\subsection{The remaining small mazes}\label{sec:cyclic}

We now account for full vertices and spatial cycles. This part of the proof
works for arbitrary initial tags. Compass directions remain absolute
throughout: translations are used to name cells, but no rotation of the fixed
controller \(\Astar\) is assumed.

\begin{lemma}[Small-maze geometry]\label{c:geometry}
Every connected induced grid maze with \(3\le |V|\le6\) belongs to exactly one
of the following structural families, with \(\mathsf X\) given priority:
\(\mathsf X\), a full vertex, its four neighbours, and at most one further
cell; \(\mathsf T\), a tree of maximum degree at most three;
\(\mathsf Q\), a unit square with at most two attached tree cells, with no
full vertex and no second square; or \(\mathsf R\), either absolute
orientation of the \(2\times3\) rectangle.
\end{lemma}
\begin{proof}
A degree-four vertex already accounts for five cells, so connectivity gives
\(\mathsf X\). In its absence an acyclic maze belongs to \(\mathsf T\).
A shortest spatial cycle in the remaining maze has length four or six,
because the grid is bipartite. A simple six-edge grid cycle has two vertical
and four horizontal edges, or four vertical and two horizontal edges. In the
first case the two vertical edges join two adjacent rows, and each row
contributes the straight two-edge segment between their endpoints. The six
vertices form a \(3\times2\) rectangle. Inducedness supplies its middle
vertical edge, and hence a unit square. Interchanging the roles of the axes
proves the same geometric assertion in the other case.

Thus a unit square is present and at most two cells remain. If these cells
create another spatial cycle, they supply a three-edge detour between two
adjacent square vertices. A one-cell exterior detour is impossible: the
common grid neighbours of opposite square vertices are already in the
square. The three-edge detour supplies a second unit square sharing an
edge with the first, and their union is a \(2\times3\) rectangle. Otherwise
the remaining cells are trees attached at single square vertices, giving
\(\mathsf Q\). The four families are disjoint under the stated priority.
\end{proof}

% Four exact baseline representatives; all grid units are 7 mm.
\begin{figure}[tbp]
\centering
\begin{tikzpicture}[maze,x=1mm,y=1mm]
 \draw[rvGray,line width=.6pt] (77,-10)--(77,51);
 \draw[rvGray,line width=.6pt] (0,23)--(158,23);
 \begin{scope}[shift={(10,31)},x=7mm,y=7mm]
  \draw[edge] (0,0)--(0,1)--(0,2); \draw[edge] (0,1)--(1,1)--(2,1);
  \foreach \x/\y in {0/0,0/1,0/2,1/1,2/1}{\node[cell] at (\x,\y) {};}
  \node[font=\footnotesize,left] at (0,1) {$J$};
 \end{scope}
 \node[paneltitle] at (33,46) {$T$\quad Subcubic trees};
 \node[note] at (33,35) {Remaining trees after\\excluding full vertices.};
 \begin{scope}[shift={(96,38)},x=7mm,y=7mm]
  \draw[edge] (-1,0)--(0,0)--(1,0)--(1,1)--(0,1)--(0,0)--(0,-1);
  \foreach \x/\y in {0/0,0/1,1/0,0/-1,-1/0,1/1}{\node[cell] at (\x,\y) {};}
  \node[font=\footnotesize,below left] at (0,0) {$p$};
 \end{scope}
 \node[paneltitle] at (116,46) {$X$\quad Full vertex};
 \node[note] at (116,35) {Takes priority,\\even with a cycle.};
 \begin{scope}[shift={(10,8)},x=7mm,y=7mm]
  \draw[edge] (0,0)--(1,0)--(1,1)--(0,1)--(0,0)--(0,-1)--(0,-2);
  \foreach \x/\y in {0/0,1/0,1/1,0/1,0/-1,0/-2}{\node[cell] at (\x,\y) {};}
  \node[font=\footnotesize,left] at (0,0) {$a$};
 \end{scope}
 \node[paneltitle] at (33,13) {$Q$\quad One square};
 \node[note] at (33,1) {At most two attached tree cells;\\no full vertex.};
 \begin{scope}[shift={(89,4)},x=7mm,y=7mm]
  \draw[edge] (0,0)--(1,0)--(2,0)--(2,1)--(1,1)--(0,1)--(0,0);
  \draw[edge] (1,0)--(1,1);
  \foreach \x/\y in {0/0,1/0,2/0,0/1,1/1,2/1}{\node[cell] at (\x,\y) {};}
 \end{scope}
 \node[paneltitle] at (116,13) {$R$\quad Rectangle};
 \node[note] at (116,2) {$2\times3$; both absolute\\orientations.};
 \begin{scope}[shift={(155,-8)},x=7mm,y=7mm]\mazeCompass{0}{0}\end{scope}
\end{tikzpicture}
\caption{The structural partition for $3\le |V|\le6$, at a common grid scale. These are geometric families, not compass symmetry classes; no start pair is selected.}
\label{f:maze-families}
\end{figure}


\begin{lemma}[Full-centre routing]\label{c:full-centre}
On every maze in \(\mathsf X\), two copies of \(\Astar\) rendezvous from
arbitrary initial tags.
\end{lemma}
\begin{proof}
Translate the full vertex to \(p=(0,0)\), and write
\(n=(0,1),e=(1,0),s=(0,-1),w=(-1,0)\). If a sixth cell \(x\) is present,
it must neighbour an arm, so
\[
x\in\{(0,2),(2,0),(0,-2),(-2,0),(1,1),(-1,1),(1,-1),(-1,-1)\}.
\]
Table~\ref{c:cross-cycles} lists all tagged cycles. A parenthesized word is
read cyclically. We verify both the words and their exhaustion directly from
Table~\ref{tab:astar}.

\begin{table}[!htbp]
\centering\small
\begin{tabular}{lll}
\toprule
Extra cell & Complete list of tagged cycles & Support type \\
\midrule
None, \((0,2),(-2,0),(-1,1)\) & \((p_1,s_1)\) & One edge \\
\((0,-2)\) & \((s_1,x_1)\) & One edge \\
\((2,0)\) & \((p_0,e_0)\), \((p_1,s_1)\) & Edges meeting at \(p\) \\
\((1,1)\) & \((p_0,e_0,x_1,e_1)\), \((p_1,s_1)\) & Shuttle and edge at \(p\) \\
\((1,-1)\) & \((p_0,e_0)\) & One edge \\
\((-1,-1)\) & \((s_1,x_0,w_0,x_1)\) & Path \(s-x-w\) \\
\bottomrule
\end{tabular}
\caption{All full-centre cases through six cells, grouped by their tagged
recurrence. Coordinates are absolute relative to the full vertex \(p\).}
\label{c:cross-cycles}
\end{table}

In the first row, \(p_0\to e_0\to p_1\leftrightarrow s_1\), while
\(e_1\to p_0\) and \(s_0\to p_0\). With no extra cell the north and west
arms send both tags to \(p_1\). For a north extension both tags at \(n\)
go to \(p_1\), and both tags at \(x\) go to \(n_1\). For a west extension,
\(w_0\to x_0\to w_1\to p_1\), and \(x_1\to w_1\). For the northwest
diagonal the remaining transitions are
\[
w_0\to x_0\to w_1\to p_1,\qquad
n_0\to x_0,\quad n_1\to p_1,\quad x_1\to n_1.
\]
These routes include every tag in those four geometries.

For a south extension both tags at \(s\) go to \(x_1\), and
\(x_1\to s_1\), \(x_0\to s_0\). The north and west leaves enter
\(p_1\); the east leaf returns to \(p_1\) from state zero and to \(p_0\)
from state one. Since \(p_0\to e_0\to p_1\to s_1\), every tag reaches
the displayed outer edge.

For an east extension the horizontal straight rule gives
\(p_0\leftrightarrow e_0\), and \(p_1\leftrightarrow s_1\).
All other tags enter those words through
\[
e_1\to x_1\to e_0,\quad x_0\to e_1,\quad
n_q\to p_1,\quad w_q\to p_1,\quad s_0\to p_0.
\]
Here and below \(q\in\{0,1\}\). For the northeast diagonal the east arm
has mask \(\N\W\), and \(x\) has mask \(\Sdir\W\), producing
\(p_0\to e_0\to x_1\to e_1\to p_0\). The southern edge remains
\((p_1,s_1)\), and the remaining tags satisfy
\[
n_0\to p_1,\quad n_1\to x_1,\quad w_q\to p_1,
\quad x_0\to n_0,\quad s_0\to p_0.
\]
In each of these two geometries one tagged cycle has edge support and the
other visits \(p\), so Lemma~\ref{t:edge-contact} applies to cross-cycle
pairs.

For the southeast diagonal the east arm has mask \(\Sdir\W\), giving
\(p_0\leftrightarrow e_0\). The other tags enter this edge through
\[
p_1\to s_1\to x_1\to s_0\to p_0,
\qquad e_1\to x_1,\qquad x_0\to e_1,
\]
together with the north and west leaf returns to \(p_1\). Finally, for the
southwest diagonal the masks of \(s,x,w\) are
\(\N\W,\N\E,\E\Sdir\). They give
\(s_1\to x_0\to w_0\to x_1\to s_1\). The remaining tags follow
\[
p_1\to s_1,\quad p_0\to e_0\to p_1,\quad e_1\to p_0,
\quad n_q\to p_1,\quad w_1\to p_1,\quad s_0\to p_1.
\]
Every single-cycle case has path support, and every tagged cycle in the
two-cycle cases also has path support. Lemma~\ref{t:same-cycle} handles
pairs entering the same tagged cycle. This proves the claim for all nine
geometric possibilities and all initial tags.
\end{proof}

\begin{lemma}[Short twig returns]\label{c:twig}
Attach a one- or two-cell path to a core vertex \(p\), with no further
neighbours at the new cells, and stop a trajectory at its first return to
\(p\). For a one-cell twig in direction \(d\), the returned state is one
if \(d\in\{\N,\W\}\), is the incoming state if \(d=\Sdir\), and is
its complement if \(d=\E\).

For a two-cell twig \(p-u-v\), write \(de\) for the directions from
\(p\) to \(u\) and from \(u\) to \(v\). All four twig tags return,
except for the following three cases: \(\E\Sdir\) has the internal edge
cycle \((u_1,v_1)\), with \(u_0,v_0\) returning to \(p_0\);
\(\Sdir\Sdir\) sends every tag to \((u_1,v_1)\); and
\(\Sdir\W\) sends every tag other than \(u_0\) to \((u_1,v_0)\),
while \(u_0\to p_1\). Every return in the other cases takes at most four
steps. The exact maps are given in Table~\ref{c:twig-table}.
\end{lemma}
\begin{proof}
The one-cell rule is the corresponding opposite-direction endpoint row of
Table~\ref{tab:astar}. For two cells there are exactly twelve words \(de\),
because \(e\ne-d\). Table~\ref{c:twig-table} prints all four successors
for each word, so following its arrows proves every stated return or
absorption. For example, the \(\E\E\) row has
\(v_0\to u_1\to v_1\to u_0\to p_0\); the \(\Sdir\Sdir\) row has
\(u_0,u_1\to v_1\to u_1\) and \(v_0\to u_0\); and the
\(\Sdir\W\) row has \(u_0\to p_1\),
\(u_1\to v_0\to u_1\), and \(v_1\to u_1\).
The appendix verifies the remaining rows and the square-port restrictions
used below.
\end{proof}

\begin{lemma}[Square routing]\label{c:square}
On every maze in \(\mathsf Q\), two copies of \(\Astar\) rendezvous from
arbitrary initial tags.
\end{lemma}
\begin{proof}
Translate the unique square so that
\[
a=(0,0),\quad b=(1,0),\quad c=(1,1),\quad d=(0,1).
\]
At most one exterior port occurs at each corner, because full vertices were
assigned to \(\mathsf X\). Table~\ref{c:square-table} gives all corner
responses. Expressions such as \((a+\Sdir)_1\) name the exterior
neighbour and its arriving state.

\begin{table}[!htbp]
\centering\small
\begin{tabular}{llll}
\toprule
Corner & Mask & Successor from state 0 & Successor from state 1 \\
\midrule
\(a\) & \(\N\E\) & \(d_0\) & \(b_1\) \\
 & \(\N\E\Sdir\) & \((a+\Sdir)_1\) & \(b_1\) \\
 & \(\N\E\W\) & \((a+\W)_0\) & \(b_1\) \\
\addlinespace
\(b\) & \(\N\W\) & \(c_1\) & \(a_0\) \\
 & \(\N\E\W\) & \(a_0\) & \((b+\E)_1\) \\
 & \(\N\Sdir\W\) & \((b+\Sdir)_0\) & \((b+\Sdir)_0\) \\
\addlinespace
\(c\) & \(\Sdir\W\) & \(d_0\) & \(b_1\) \\
 & \(\N\Sdir\W\) & \(b_0\) & \(b_0\) \\
 & \(\E\Sdir\W\) & \(b_1\) & \(b_1\) \\
\addlinespace
\(d\) & \(\E\Sdir\) & \(a_1\) & \(c_1\) \\
 & \(\N\E\Sdir\) & \(a_1\) & \(c_1\) \\
 & \(\E\Sdir\W\) & \(a_1\) & \(a_1\) \\
\bottomrule
\end{tabular}
\caption{Every possible nonfull square-corner row. Exterior ports occur in
at most two additional cells in total.}
\label{c:square-table}
\end{table}

The exterior ports at \(c\) and \(d\) are never selected. This fact alone
does not exclude an external attractor. At ports \(c\N,c\E,d\N,d\W\),
the allowed two-cell words are, respectively,
\[
\{\N\N,\N\E\},\quad \{\E\E,\E\N\},\quad
\{\N\N,\N\W\},\quad \{\W\W,\W\N\}.
\]
The omitted continuation would touch the neighbouring square corner and
produce \(\mathsf R\). Lemma~\ref{c:twig} proves that every tag in every
allowed dormant twig returns to the square. Once such a tag has returned,
the twig cannot be entered again. Also \(d\) can be entered recurrently
only from \(a\) by a northward move arriving in state zero, or from
\(c\) by a westward move arriving in state zero. Hence \(d_1\) is
transient, even though a dormant-twig root may visit it once.

The active ports are \(a\Sdir,a\W,b\E,b\Sdir\). Their complete
return rules are in Table~\ref{c:active-ports}. If one port already uses a
cell, each other twig has at most one cell. Every tag on a one-cell twig
returns immediately to the square; the specific return states used below
refer to its entry from the core. The possible masks at \(b\) now give
four routing regimes.

\emph{First, let \(b=\N\W\) and
\(c\in\{\Sdir\W,\E\Sdir\W\}\).}
Both tags at \(c\) feed the left core, whose fixed route is
\[
d_0\longrightarrow a_1\longrightarrow b_1\longrightarrow a_0.
\]
If \(a=\N\E\), it closes through \(a_0\to d_0\), giving the unique
tagged cycle \((d_0,a_1,b_1,a_0)\), with path support \(d-a-b\).
If \(a\) has a south or west port, only \(a_0\) leaves through it.
A one-cell twig returns to \(a_1\), as does every tag on a two-cell west
twig. The resulting unique tagged cycle is a shuttle extending \(b-a\)
by that twig. A two-cell south twig has word \(\Sdir\Sdir\) or
\(\Sdir\W\), and \(a_0\) enters it as \(u_1\), which belongs to
its absorbing basin. In the latter case the only returning tag is
\(u_0\), and it follows
\(u_0\to a_1\to b_1\to a_0\to u_1\).
Thus all roots reach the same internal edge cycle. The unused tag
\(b_0\) enters \(c_1\), and all other exterior tags return from dormant
twigs, so no other tagged cycle remains.

\emph{Second, let \(b=\N\W\) and \(c=\N\Sdir\W\).}
The north port at \(c\) costs at least one exterior cell. There is an edge
cycle \(b_0\to c_1\to b_0\). The left route
\(d_0\to a_1\to b_1\to a_0\) closes through \(d_0\), or through
a one-cell twig at \(a\) returning to \(a_1\). That twig cannot have
two cells because of the port at \(c\). Hence the other tagged cycle is a
four-tag shuttle on a three-vertex path and contains \(b_1\). The
dormant twig at \(c\) returns and then enters \(b_0\); a dormant twig at
\(d\) also returns to the square. The displayed corner transitions
exhaust all other tags. The two supports share \(b\), and one is an edge,
so Lemma~\ref{t:edge-contact} applies.

\emph{Third, let \(b=\N\E\W\).}
There is an eastern twig, \(b_0\to a_0\), and \(b_1\) exits east in
state one. Every square tag eventually reaches \(b_1\): the tags at
\(c\) enter \(b_0,b_1\), or \(d_0\), while \(d_0\to a_1\to b_1\).
At \(a_0\), either the route through \(d_0\) reaches \(b_1\), or a
twig at \(a\) returns to \(a_1\) and then \(b_1\). Such a twig has
at most one cell. A one-cell eastern twig returns to \(b_0\), and a
two-cell \(\E\E\) twig returns every tag to \(b_0\). In these cases
there is one tagged cycle, supported on the path from the left cap
(\(d\), or the single twig cell at \(a\)) through \(a,b\) to the
eastern twig. If the eastern twig has word \(\E\Sdir\), the departure
from \(b_1\) enters its edge absorber. Its returnable tags \(u_0,v_0\)
first return to \(b_0\), then reach \(b_1\), and next enter the same
absorber. Thus this alternative has a unique edge-supported tagged cycle.
Dormant twig returns again account for all exterior roots.

\emph{Fourth, let \(b=\N\Sdir\W\).}
Both tags at \(b\) exit south in state zero. Every square tag reaches
\(b\); any twig at \(a\) has at most one cell and returns to \(a_1\),
which goes directly to \(b_1\). A one-cell southern twig returns to
\(b_0\), giving an edge cycle. A two-cell \(\Sdir\E\) twig also
returns all tags to \(b_0\), and the tagged cycle still uses only \(b\)
and its first twig cell. A two-cell \(\Sdir\Sdir\) twig captures every
tag on its outer state-one edge. All other roots feed the indicated unique
tagged cycle.

These regimes cover all possible corner masks. All displayed tagged cycles
have path support. Lemma~\ref{t:same-cycle} handles same-cycle pairs, and
Lemma~\ref{t:edge-contact} additionally handles the second regime.
Hence all initial tags are covered.
\end{proof}

\begin{lemma}[The two rectangle orientations]\label{c:rectangles}
On both mazes in \(\mathsf R\), two copies of \(\Astar\) rendezvous from
arbitrary initial tags.
\end{lemma}
\begin{proof}
For the horizontal rectangle label the bottom row
\(a=(0,0),b=(1,0),c=(2,0)\) and the top row
\(d=(0,1),e=(1,1),f=(2,1)\).
For the vertical rectangle use
\(a=(0,0),b=(1,0),c=(0,1),d=(1,1),e=(0,2),f=(1,2)\).
The two complete maps are printed separately in
Table~\ref{c:rectangle-maps}.

\begin{table}[!htbp]
\centering\small
\begin{minipage}[t]{.48\textwidth}
\centering
Horizontal rectangle\par\smallskip
\begin{tabular}{llll}
\toprule
Cell & Mask & From 0 & From 1 \\
\midrule
\(a\)&\(\N\E\)&\(d_0\)&\(b_1\)\\
\(b\)&\(\N\E\W\)&\(a_0\)&\(c_1\)\\
\(c\)&\(\N\W\)&\(f_1\)&\(b_0\)\\
\(d\)&\(\E\Sdir\)&\(a_1\)&\(e_1\)\\
\(e\)&\(\E\Sdir\W\)&\(b_1\)&\(b_1\)\\
\(f\)&\(\Sdir\W\)&\(e_0\)&\(c_1\)\\
\bottomrule
\end{tabular}
\end{minipage}\hfill
\begin{minipage}[t]{.48\textwidth}
\centering
Vertical rectangle\par\smallskip
\begin{tabular}{llll}
\toprule
Cell & Mask & From 0 & From 1 \\
\midrule
\(a\)&\(\N\E\)&\(c_0\)&\(b_1\)\\
\(b\)&\(\N\W\)&\(d_1\)&\(a_0\)\\
\(c\)&\(\N\E\Sdir\)&\(a_1\)&\(d_1\)\\
\(d\)&\(\N\Sdir\W\)&\(b_0\)&\(b_0\)\\
\(e\)&\(\E\Sdir\)&\(c_1\)&\(f_1\)\\
\(f\)&\(\Sdir\W\)&\(e_0\)&\(d_1\)\\
\bottomrule
\end{tabular}
\end{minipage}
\caption{All twelve tagged-state transitions in each absolute rectangle
orientation. These are two separate controller calculations.}
\label{c:rectangle-maps}
\end{table}

The horizontal map has the tagged cycle
\((a_0,d_0,a_1,b_1,c_1,b_0)\), supported on \(d-a-b-c\).
Every remaining tag reaches it: \(c_0\to f_1\to c_1\),
\(d_1\to e_1\to b_1\), \(e_0\to b_1\), and
\(f_0\to e_0\). Lemma~\ref{t:same-cycle} proves contact.

The vertical map has the two tagged cycles
\((a_0,c_0,a_1,b_1)\) and \((b_0,d_1)\), supported on \(c-a-b\)
and \(b-d\). The remaining tags satisfy
\(c_1,d_0\to d_1,b_0\), respectively,
\(e_0\to c_1\), \(e_1\to f_1\to d_1\), and \(f_0\to e_0\).
Thus these words exhaust recurrence. Lemma~\ref{t:same-cycle} handles
same-cycle pairs, and Lemma~\ref{t:edge-contact} handles cross-cycle pairs
at the common vertex \(b\).
\end{proof}

\begin{proposition}[Six-cell success]\label{c:six-cell}
The controller \(\Astar\) succeeds from every pair of state-zero starts
on every connected induced grid maze with at most six cells.
\end{proposition}
\begin{proof}
On one or two cells every pair is already successful at time zero.
For larger mazes Lemma~\ref{c:geometry} is exhaustive.
The family \(\mathsf T\) is covered by Proposition~\ref{t:small-trees};
the other three families are covered by Lemmas~\ref{c:full-centre},
\ref{c:square}, and~\ref{c:rectangles}.
\end{proof}
